% COS'È QUESTO FILE?
% In questo documento, riporto le versioni vecchi di codice, estratte della classe HKNdocument03, che non sono più utili
% e che è stato di conseguenza eliminato. Le Lascio tuttavia come esempio didattico, perchè è importante anche capire
% il perchè certe soluzioni non funzionino, o siano altamente sconvenienti.

% NOTA: ovviamente buona parte del codice sarà costituito dalle parti di codice eliminate da HKNdocument02.cls, per i
% problemi spiegati nelle ultime modifiche di quella versione, in particolare per il "BUG 6", ossia la necessità di
% revisionare la logica di modifica e gestione delle testatine. È il motivo per cui sono passato alla versione 03.
